\subsection{Simulation Model for Poking}\label{sec:sim-model}

In order to understand the effects of impulsive forces on objects, we use the PyBullet physics simulation engine \cite{coumans2019}. PyBullet models rigid body dynamics by performing numerical integration over time with equations of motion to solve for object position and velocity. Joint constraints, contact forces, and friction, in addition to external forces such as gravity, are taken into account in the forward dynamics solver of the engine and modeled as constraints in a Linear Complementarity Problem (LCP).

The main point of interest when modeling impulsive interactions is contact force. In the context of poking,  contact is primarily dominated by frictional interactions between i) the robot end-effector and the object of interest, and ii) the object and the environment (i.e., the object's planar support surface and surrounding obstacles). Analytical models for contact make simplifying assumptions that do not fully represent the complexity of real-world dynamics, including the nonlinear nature of friction and actuator degradation and latency. Additionally, they do not generalize well to a diverse set of objects.
Conversely, learning-based approaches, while capable of achieving generalization and modeling uncertainty and complexity, are data-inefficient. Collecting abundant and task-representative data in real-world robotics is expensive with regard to the time taken and the wear-and-tear caused on a real robot through repetitive, and potentially, high-acceleration trajectory executions such as those characteristic to poking manipulation. Executing such trajectories and modeling collisions along the action path with the object and the environment is therefore far safer in simulation.

While simulation does not perfectly capture the inherent complexity and stochasticity of real-world contact dynamics due to the simplistic nature of the underlying analytical models used, it nonetheless provides a good balance between pure learning and analytical models by ensuring interaction modeling is both cheap and safe while encapsulating the essential characteristics of robot--object and object--environment interactions. The closed-loop nature of our proposed motion planner also compensates for any inaccuracies in simulation modeling while executing poke plans in the real-world. That is, if a resultant pose violates a predefined object pose threshold, we compute a new poke plan from the current object pose to the goal region. This crucial feature allows our planner to plan feasible poke paths in simulation and execute them in the real-world.\vspace{-4pt}
